% Chapter 3, section 7: Rate limiting
% Write in plain LaTeX. Code of a guide step: \guide{<step>}; commands: \begin{shell}...\end{shell};
% a screenshot: \screenshot{file}{caption}; a gap in the draft: \placeholder{...}.
% See book/latex/queen.sty for the full list.
%
% Draft by Claude, 2026-10-24. The code is cut from examples/tutorials/js/book-use-cases.mjs, which
% checks every result quoted here against a 2.7.0 node.

\section{Rate limiting}

An SMS provider accepts five messages a second for each customer. Send more and it refuses them,
or blocks the account. The messages must still all go out, in order.

A stream with a gate decides, message by message, whether a message may go on now. The token
bucket gate keeps its state per partition, in the broker: here five tokens per customer, refilled
at five a second. A message the gate refuses stops its partition there, and comes back, with the
ones after it, when its lease ends.

\guide{uc-rate}

Five messages went at once, and the rest followed at about five a second, so the twelfth left
about two seconds after the first: all twelve, in order.

\queenfig{sms-gate}{The token bucket gate: 5, 5, then 2, in order. A refill a few milliseconds short lets 4 through, not 5.}

\why{the bucket lives in the broker, per customer, and survives restarts: no Redis next to it.
Waiting messages stay in their queue, in order, and never need a queue of your own.}

One thing to know: the refused messages wait for their lease, so size the two together. The most a
gate can let through is the smaller of \texttt{refillPerSec} and \texttt{capacity} divided by the
lease in seconds, so keep \texttt{capacity} near \texttt{refillPerSec} times the lease.
